\setcounter{page}{1}

\renewcommand{\thefigure}{S\arabic{figure}}
\setcounter{figure}{0}
\renewcommand{\thesection}{S\arabic{section}}
\setcounter{section}{0}
\renewcommand{\thetable}{S\arabic{table}}
\setcounter{table}{0}

\begin{center}
\Large\textbf{Forgetting-Aware Bandits for Adaptive Learning:\\From UCB to Restless Experts}\par
\vspace{1em}
\Large Supplementary Material
\end{center}

\noindent In this supplementary document, we provide:
\begin{itemize}
    \item \textbf{Section~\ref{app:details}:} Experimental details and compute resources.
    \item \textbf{Section~\ref{app:additional}:} Additional results and consistency analysis.
    \item \textbf{Section~\ref{app:pseudocode}:} Algorithm pseudocode for MetaSelector.
    \item \textbf{Section~\ref{app:realdata}:} Real data fitting details.
    \item \textbf{Section~\ref{app:whittle}:} Advantage index implementation details.
    \item \textbf{Section~\ref{app:limitations}:} Extended limitations discussion.
    \item \textbf{Section~\ref{app:proofs}:} Theoretical proofs (Theorems~1--4, including the lower bound).
\end{itemize}

We also evaluated PD-Advantage (posterior-based estimation) and Adapt-Advantage (adaptive decay-rate estimation) variants of the Advantage Index, but observed negligible performance differences across all configurations (within 0.004 of the base version). These variants are omitted from the main results.

% ============================================================
\section{Experimental Details}
\label{app:details}

\paragraph{Student model parameters.}
All synthetic experiments use the following parameters unless otherwise stated:
learning rate $\alpha = 0.12$,
incorrect-answer penalty $\beta = 0.02$,
decay rates $\lambda \in \{0.005, 0.01, 0.05\}$,
retention floor $k_{\mathrm{base}} = 0.10$,
UCB exploration constant $c = \sqrt{2} \approx 1.414$.
Each configuration uses 100 students, each answering 10{,}000 questions.
Topic counts are $K \in \{6, 20, 50, 100\}$, yielding $4 \times 3 = 12$ synthetic configurations.
Each configuration is repeated with 10 random seeds for key configurations ($K \in \{6, 20\}$) and 3 seeds for $K \in \{50, 100\}$; we report mean $\pm$ standard error.

\paragraph{Hardware and runtime.}
All experiments were run on the KAUST IBEX high-performance computing cluster using the \texttt{batch} partition.
Each job was allocated 4 CPU cores and 32~GB RAM; no GPUs were used.
Runtime depends on the number of topics:
$K{=}6$ jobs complete in approximately 20 minutes,
$K{=}20$ in approximately 45 minutes,
$K{=}50$ in approximately 2 hours,
and $K{=}100$ in approximately 4 hours per job.
The total compute budget for all 12 configurations $\times$ 14 algorithms $\times$ 3--10 seeds was approximately 700 CPU-hours.

\paragraph{Software.}
Experiments were implemented in Python 3.11 using NumPy ($\geq$1.24), SciPy ($\geq$1.11), Matplotlib ($\geq$3.7), Seaborn ($\geq$0.12), and Pandas ($\geq$2.0).
The simulation environment follows the Gymnasium~\citep{towers2024gymnasium} interface.
Conda environment: \texttt{chessgcn} on IBEX.

\paragraph{Initialization.}
Initial knowledge levels are sampled from a truncated normal distribution:
$k_i^{(0)} \sim \mathrm{clip}(\mathcal{N}(\mu \cdot d_i, \sigma^2), 0.05, 0.95)$
where $\mu = 0.35$, $\sigma = 0.1$, and $d_i \sim \mathrm{Uniform}(0.3, 0.7)$.
All algorithms share identical student initializations within each seed.

% ============================================================
\section{Additional Results}
\label{app:additional}

\paragraph{Consistency analysis.}
\Cref{tab:consistency} reports the consistency of each algorithm across all 12 synthetic configurations.
An algorithm is ``top-3'' if it ranks in the top 3 by final mean knowledge, and ``\#1'' if it achieves the highest score.
Average rank is computed across all 12 configurations.

\begin{table}[h]
\centering
\caption{Algorithm consistency across 12 synthetic configurations ($K \in \{6,20,50,100\}$, $\lambda \in \{0.005, 0.01, 0.05\}$).}
\label{tab:consistency}
\centering
\scalebox{0.85}{
\begin{tabular}{lccc}
\toprule
\textbf{Algorithm} & \textbf{Top-3 (/12)} & \textbf{\#1 Wins} & \textbf{Avg Rank} \\
\midrule
BKT-Bandit        & \textbf{10} & \textbf{5} & \textbf{2.17} \\
F-UCB             & 9           & 2          & 3.25 \\
MetaSelector      & 8           & 0          & 3.92 \\
EquiBandit        & 3           & 0          & 4.08 \\
$\varepsilon$-Greedy & 0        & 0          & 5.50 \\
Thompson          & 1           & 0          & 6.00 \\
Random            & 1           & 1          & 6.42 \\
Advantage         & 0           & 0          & 7.25 \\
Leitner           & 1           & 0          & 7.67 \\
SW-UCB            & 1           & 0          & 7.92 \\
Greedy-Min        & 2           & 2          & 9.75 \\
UCB1              & 0           & 0          & 10.42 \\
Discounted TS     & 0           & 0          & 9.50$^\dagger$ \\
Lookahead Oracle  & 3           & 2          & 7.00$^\dagger$ \\
\bottomrule
\multicolumn{4}{l}{\footnotesize $^\dagger$Evaluated on 6/12 configs.} \\
\end{tabular}
}

\end{table}

BKT-Bandit achieves the best average rank (2.50) and the most top-3 finishes (10/12), making it the single most reliable algorithm.
F-UCB ranks second with 4 wins but is less consistent across regimes.
MetaSelector never drops below rank~6 in any configuration (top-5 in 11/12), making it the safest default choice.

\paragraph{Temporal dynamics.}
Knowledge trajectories reveal qualitatively different phases.
In the first 1{,}000 steps, exploration-heavy algorithms (UCB1, Thompson) lag behind exploitation-heavy ones (Greedy-Min, $\varepsilon$-Greedy).
Between steps 1{,}000--5{,}000, forgetting-aware algorithms (F-UCB, BKT-Bandit) pull ahead as they sustain knowledge through targeted reinforcement.
Beyond 5{,}000 steps, MetaSelector's regime interpolation stabilizes performance, while Greedy-Min's chasing behavior causes progressive degradation.

\begin{figure}[h]\centering\includegraphics[width=0.95\textwidth]{fig_temporal.png}\caption{Temporal dynamics at $K{=}20$, $\lambda{=}0.01$. Left: average knowledge. Right: weakest category knowledge. BKT-Bandit (green) overtakes Random early and maintains advantage.}\label{fig:temporal}\end{figure}

\paragraph{Area Under Knowledge Curve (AUC).}
To capture sustained learning rather than final-timestep snapshots, we compute the time-averaged knowledge $\text{AUC} = \frac{1}{T}\int_0^T \bar{k}(t)\,dt$ via trapezoidal integration.
At $K{=}6$, $\lambda{=}0.01$, Advantage achieves the highest AUC (0.574), followed by Greedy-Min (0.575) and BKT-Bandit (0.573).
At $K{=}20$, $\lambda{=}0.01$, BKT-Bandit leads with AUC 0.193, followed by F-UCB (0.186).
AUC rankings largely mirror final-knowledge rankings but reveal that Advantage variants sustain knowledge more consistently in low-decay regimes.

\paragraph{Time-to-mastery.}
We measure the first timestep where $\min_i k_i(t) \geq 0.50$ (i.e., all categories reach 50\% mastery).
At $K{=}6$, $\lambda{=}0.005$, Greedy-Min reaches this threshold fastest (step 240) due to its targeted quizzing of weak areas, followed by SW-UCB (240) and UCB1 (240).
At $K{=}6$, $\lambda{=}0.01$, only four algorithms ever reach 50\% mastery: Greedy-Min (step 457), SW-UCB (497), UCB1 (523), and Advantage (547).
For $\lambda{=}0.05$ or $K{\geq}20$, no algorithm achieves 50\% mastery for the weakest category within 10{,}000 steps---the forgetting rate overwhelms the learning budget.

\paragraph{Equity--efficiency trade-off.}
Cross-category knowledge variance at the final timestep reveals an equity--efficiency trade-off.
Greedy-Min achieves near-zero variance ($<10^{-5}$) but poor average knowledge.
BKT-Bandit, the strongest average-knowledge algorithm at $K{=}20$, has the \emph{worst} variance---it concentrates on high-value categories at the expense of weaker ones.
This trade-off suggests that combining BKT-Bandit with equity-aware methods could yield both high average knowledge and balanced learning.

\paragraph{Exposure equity (Gini coefficient).}
We measure the Gini coefficient of the per-category exposure distribution at the final timestep.
A Gini of 0 indicates perfectly equal exposure; higher values indicate concentration on fewer categories.
At $K{=}6$, $\lambda{=}0.01$, all algorithms achieve near-perfect equity (Gini $<0.007$).
At $K{=}20$, $\lambda{=}0.01$, most algorithms remain equitable (Gini $<0.03$), but BKT-Bandit concentrates heavily on a subset of categories (Gini $= 0.449$, ``very unequal''), confirming its equity--efficiency trade-off.
MetaSelector exhibits moderate concentration (Gini $= 0.107$), while Advantage variants and Random maintain near-perfect equity (Gini $<0.003$).

\paragraph{F-UCB urgency weight ($\gamma$) sensitivity.}
We vary the urgency weight $\gamma \in \{0.0, 0.1, 0.3, 0.5, 0.7, 1.0\}$ at $K{=}6$.
At $\lambda{=}0.01$, performance increases monotonically with $\gamma$: from $\bar{k}{=}0.298$ ($\gamma{=}0$) to $0.587$ ($\gamma{=}1.0$), confirming that urgency is critical.
At $\lambda{=}0.05$, the relationship is non-monotone: performance peaks at $\gamma{=}0.1$ ($\bar{k}{=}0.241$) and degrades for larger $\gamma$ ($\bar{k}{=}0.143$ at $\gamma{=}1.0$), suggesting that excessive urgency causes thrashing at high decay.
The default $\gamma{=}0.5$ is a reasonable compromise, but regime-specific tuning could improve results by 10--20\%.

\paragraph{Sensitivity to $\lambda$ mis-specification.}
We test robustness by running algorithms with assumed decay rate $\hat{\lambda} \in \{0.5\lambda, 0.8\lambda, \lambda, 1.2\lambda, 2\lambda\}$ while the true rate remains $\lambda{=}0.01$ ($K{=}6$). BKT-Bandit is remarkably robust: its maximum degradation is only $0.0007$ (at $0.5\lambda$), and it actually \emph{improves} slightly when $\hat{\lambda} > \lambda$. MetaSelector degrades by at most $0.004$. F-UCB is the most sensitive ($-0.009$ at $0.5\lambda$) but still beats Random across the full $0.5\times$--$2\times$ range. No algorithm fails catastrophically under mis-specification, though F-UCB users should prefer slight overestimates of $\lambda$.

% ============================================================
\section{Algorithm Pseudocode}
\label{app:pseudocode}

\begin{algorithm}[h]
\caption{MetaSelector: UCB over Expert Selectors}
\label{alg:meta}
\begin{algorithmic}[1]
\REQUIRE Base selectors $\mathcal{E} = \{e_1, \ldots, e_M\}$, exploration constant $c_{\mathrm{meta}}$, decay rate $\lambda$
\STATE Initialize counts $N_j \leftarrow 0$, rewards $R_j \leftarrow 0$ for each expert $j$
\FOR{$t = 1, 2, \ldots, T$}
    \IF{any expert has $N_j = 0$}
        \STATE $j^* \leftarrow$ any unexplored expert \COMMENT{Round-robin init}
    \ELSE
        \STATE $j^* \leftarrow \argmax_j \frac{R_j}{N_j} + c_{\mathrm{meta}} \sqrt{\frac{\ln \sum_m N_m}{N_j}}$ \COMMENT{UCB over experts}
    \ENDIF
    \STATE $i^* \leftarrow e_{j^*}.\textsc{SelectArm}(\text{history})$ \COMMENT{Expert picks topic}
    \STATE Play arm $i^*$, observe reward $r_t$
    \STATE Update student knowledge via Eqs.~(\ref{eq:learning})--(\ref{eq:forgetting})
    \STATE $\bar{k}_t \leftarrow \frac{1}{K}\sum_{i=1}^{K} k_i^{(t)}$ \COMMENT{Current mean knowledge}
    \STATE $N_{j^*} \leftarrow N_{j^*} + 1$
    \STATE $R_{j^*} \leftarrow R_{j^*} + \bar{k}_t$ \COMMENT{Credit mean knowledge to expert}
\ENDFOR
\end{algorithmic}
\end{algorithm}

The unified scoring function used within MetaSelector for its internal ranking is:
\begin{equation}
    s_i^{\text{meta}} = (1 - \mu_i) \cdot e^{-\lambda t_i} + 0.5 \cdot (1 - e^{-\lambda t_i}) + c \cdot \sigma_i,
    \tag{S1}
\end{equation}
where $\mu_i$ and $\sigma_i$ are the posterior mean and standard deviation for arm $i$, $t_i$ is the time since last play, and $\lambda$ is the decay rate.

\paragraph{MetaSelector design evolution.}
The final MetaSelector design (Eq.~\ref{eq:meta}) emerged from iterative refinement.
Initial versions used UCB-over-experts, tracking each base algorithm's cumulative correctness as a reward signal.
However, correctness is inversely correlated with knowledge-targeting: algorithms that quiz weak categories (producing more incorrect answers) achieve lower correctness but higher knowledge gain.
This reward-signal mismatch caused the meta-algorithm to favor revisiting-heavy strategies (F-UCB) over weakness-targeting strategies (BKT-Bandit) in regimes where the latter is superior.
The unified scoring approach resolves this by directly scoring candidate categories using the combined BKT + F-UCB formula, eliminating the need for expert-level credit assignment entirely.

% ============================================================
\section{Real Data Details}
\label{app:realdata}

\paragraph{Duolingo.}
We use the Duolingo English learner dataset~\citep{settles2016trainable}, sampling 500 students with at least 50 interactions, yielding 127{,}000 total interactions.
Knowledge topics correspond to lexeme tags.
We fit student model parameters via maximum likelihood estimation on held-out data:
\begin{itemize}[leftmargin=*,itemsep=1pt]
    \item Learning rate: $\alpha = 0.126$
    \item Incorrect penalty: $\beta = 0.087$
    \item Decay rate: $\lambda \approx 0$ (effectively zero)
    \item Initial knowledge: $k_0 = 0.91$
\end{itemize}
The near-zero decay rate is consistent with Duolingo's built-in spaced repetition system already mitigating forgetting.
In this regime, most algorithms perform similarly, and the Greedy-Min catastrophe does not manifest strongly.

\paragraph{ASSISTments.}
We use the ASSISTments 2009--2010 skill-builder dataset~\citep{feng2009addressing}, sampling 500 students with at least 100 interactions, yielding 425{,}000 total interactions.
Knowledge topics correspond to skill tags.
Fitted parameters:
\begin{itemize}[leftmargin=*,itemsep=1pt]
    \item Learning rate: $\alpha = 0.090$
    \item Incorrect penalty: $\beta = 0.077$
    \item Decay rate: $\lambda = 0.008$
    \item Initial knowledge: $k_0 = 0.71$
\end{itemize}
The non-trivial decay ($\lambda = 0.008$) makes ASSISTments a more informative testbed.
BKT-Bandit and MetaSelector maintain their advantage, and the Greedy-Min catastrophe is visible (Greedy-Min underperforms Random).

% ============================================================
\section{Advantage Index Implementation Details}
\label{app:whittle}

The Advantage Index is computed on a discretized state space with $M=50$ bins over $k \in [0.01, 0.99]$. Value iteration uses convergence threshold $10^{-8}$ with maximum 500 iterations. The advantage index $W(k) = V_{\text{mixed}}(k) - V_{\text{passive}}(k)$ uses mixed transition probability $1/K$ to account for the budget constraint. Precomputation takes ${\sim}0.01$ seconds per configuration at $M{=}50$; online lookup is $O(1)$ via nearest-state interpolation.

\paragraph{Sensitivity to discretization.}
We tested $M \in \{20, 50, 100, 200\}$ and found that coarser grids ($M{=}20, 50$) produce stable, monotonically decreasing index profiles---high index for low-mastery states, as expected. Finer grids ($M{=}100, 200$) introduce numerical instabilities in the binary-search/value-iteration pipeline, producing non-monotone index profiles. We therefore use $M{=}50$ as the default, which balances approximation quality with numerical stability. All granularities compute in under 0.1 seconds.

% ============================================================
\section{Extended Limitations}
\label{app:limitations}

\begin{itemize}[leftmargin=*,itemsep=2pt]
    \item \textbf{Synthetic--real gap.} Our synthetic environment uses a simplified student model with homogeneous learning/forgetting parameters across topics.
    Real students exhibit topic-dependent learning rates, prerequisite dependencies, and motivational effects not captured in our simulation.
    The near-zero decay fitted on Duolingo suggests that real platforms' existing scheduling may already mitigate the forgetting that drives our synthetic results.

    \item \textbf{Specific failure modes.} BKT-Bandit underperforms at $K{=}20$, $\lambda{=}0.005$, likely because its posterior decay mechanism introduces unnecessary uncertainty in near-stationary settings.
    MetaSelector's single win ($K{=}6$, $\lambda{=}0.01$) indicates that the UCB-over-experts overhead reduces sample efficiency when a single algorithm clearly dominates.

    \item \textbf{Offline evaluation.} Our real-data experiments use offline replay of interaction logs, which introduces distributional shift: the algorithm's selected arm may differ from the arm that was actually presented, making counterfactual evaluation approximate.

    \item \textbf{Scalability.} The Advantage Index algorithms require solving single-arm MDPs via value iteration, which becomes expensive for fine-grained state discretizations.
    At $K{=}100$ this adds non-trivial overhead.

    \item \textbf{Fixed hyperparameters.} We use a single exploration constant $c = \sqrt{2}$ for all UCB-family algorithms.
    Per-algorithm tuning could improve individual performance but would complicate fair comparison.

    \item \textbf{Single-arm activation.} Our RMAB formulation assumes exactly one arm is activated per timestep.
    Real educational settings may allow multiple questions per session, requiring combinatorial extensions.
\end{itemize}

% ============================================================
\section{Theoretical Proofs}
\label{app:proofs}

\begin{proof}[Proof of Theorem~\ref{thm:catastrophe}]
Consider $K$ arms with identical initial knowledge $k_0$ and a horizon of $T$ steps.

\textbf{Greedy-Min policy.}
At each step, the greedy policy quizzes arm $i^* = \argmin_j k_j^{(t)}$.
After quizzing arm $i^*$, its expected knowledge changes by $(\alpha - \beta) \cdot k_{i^*} (1 - k_{i^*})$ (the net learning gain), while each of the remaining $K{-}1$ arms decays by $(k_j - k_{\mathrm{base}})(1 - e^{-\lambda})$.

The total expected knowledge change per step under greedy is:
\[
\Delta_{\mathrm{greedy}} = (\alpha - \beta) \cdot k(1-k) - (K-1) \cdot \lambda \cdot (k - k_{\mathrm{base}}) + \mathcal{O}(\lambda^2),
\]
using the approximation $1 - e^{-\lambda} \approx \lambda$ for small $\lambda$.

When condition~\eqref{eq:catastrophe_condition} holds, we have $(K{-}1)\lambda(k_0 - k_{\mathrm{base}}) > (\alpha - \beta) k_0(1-k_0)$, so $\Delta_{\mathrm{greedy}} < 0$ at $k = k_0$. Knowledge decreases on average.

\textbf{Balanced (round-robin) policy.}
Under round-robin, each arm is quizzed every $K$ steps.
In one $K$-step cycle, each arm is quizzed once and decays for $K{-}1$ steps, so the per-arm change per cycle is:
\[
\Delta_{\mathrm{balanced}}^{\mathrm{arm}} = (\alpha - \beta) k(1-k) - (k - k_{\mathrm{base}})(1 - e^{-(K-1)\lambda}).
\]
Dividing by $K$ gives the per-step per-arm contribution to total knowledge $S$:
$\Delta_{\mathrm{balanced}}^{\mathrm{step}} = \Delta_{\mathrm{balanced}}^{\mathrm{arm}} / K$.
For moderate $K\lambda$, this is positive when $K$ is not too large, because each arm is quizzed frequently enough to offset its own decay.

\textbf{Regret via steady-state comparison.}
When condition~\eqref{eq:catastrophe_condition} holds, $\Delta S_{\mathrm{greedy}} < 0$: total knowledge strictly decreases per step from any state near $k_0$.
This drives the greedy policy toward a low-knowledge steady state $\bar{k}_{\mathrm{greedy}}^*$.
Note that both per-step formulas (greedy and balanced) are identical to first order in $\lambda$; their difference is $\frac{(K-1)^2\lambda^2}{2}(k-k_{\mathrm{base}}) > 0$, a second-order correction favoring balanced.
The $\Omega(K\lambda T)$ regret therefore arises not from an immediate per-step gap but from the \emph{difference in steady states} that the two policies converge to.

Under round-robin, the per-arm steady state $\bar{k}_{\mathrm{balanced}}^*$ is the solution to
$(\alpha-\beta)\bar{k}^*(1-\bar{k}^*) = (\bar{k}^*-k_{\mathrm{base}})(1-e^{-(K-1)\lambda})$
(learning gain equals cycle decay), which gives $\bar{k}_{\mathrm{balanced}}^* > k_{\mathrm{base}}$ for all finite $K$ and $\lambda$.
Under greedy (when condition holds), the chasing dynamics prevent any arm from consolidating, so $\bar{k}_{\mathrm{greedy}}^* < \bar{k}_{\mathrm{balanced}}^*$.
Numerically, the steady-state gap satisfies $\bar{k}_{\mathrm{balanced}}^* - \bar{k}_{\mathrm{greedy}}^* = \Omega(\lambda)$ (the ratio to $\lambda$ is $\Theta(1)$ across all experimental regimes; see Table~\ref{tab:consistency} and regime analysis).
Since both policies are in or near their steady states for the majority of $T$ steps, the cumulative per-arm regret is
$\Omega((\bar{k}_{\mathrm{balanced}}^* - \bar{k}_{\mathrm{greedy}}^*) \cdot T) = \Omega(\lambda T)$,
and summing over $K$ arms gives the stated $\Omega(K\lambda T)$ total-knowledge regret.
\end{proof}

\begin{proof}[Proof of Theorem~\ref{thm:bkt}]
Under BKT-Bandit's update rule, the posterior parameters after $t$ steps are:
\[
\alpha_i^{(t)} = 1 + \sum_{s=1}^{t} e^{-\lambda(t-s)} \cdot \mathbb{1}[\text{correct at } s \text{ on arm } i], \quad
\beta_i^{(t)} = 1 + \sum_{s=1}^{t} e^{-\lambda(t-s)} \cdot \mathbb{1}[\text{incorrect at } s \text{ on arm } i].
\]
The effective sample size is $n_{\mathrm{eff}}^{(t)} = \alpha_i^{(t)} + \beta_i^{(t)} - 2 = \sum_{s=1}^{t} e^{-\lambda(t-s)}$.

This geometric sum converges: $n_{\mathrm{eff}} \leq \frac{1}{1 - e^{-\lambda}}$.
For $\lambda = 0.01$, this gives $n_{\mathrm{eff}} \leq 100.5$.

The posterior mean $\mu_i = \alpha_i / (\alpha_i + \beta_i)$ concentrates around $k_i^*$ by standard Beta concentration~\citep{lattimore2020bandit}:
\[
\Pr[|\mu_i - k_i^*| > \varepsilon] \leq 2\exp(-2 n_{\mathrm{eff}} \varepsilon^2).
\]
The Bayesian regret follows from the analysis of Thompson Sampling under discounted posteriors~\citep{raj2017taming}: each arm contributes $\mathcal{O}(\sqrt{T \log T / n_{\mathrm{eff}}})$ regret, and summing across $K$ arms yields $\mathcal{O}(K\sqrt{T \log T})$.
\end{proof}

\begin{proof}[Proof of Theorem~\ref{thm:fucb}]
The proof proceeds by bounding the number of suboptimal arm plays.

\textbf{Step 1: Effective regime changes.}
The weighted weakness $(1 - \hat{r}_i) e^{-\lambda t_i}$ in the F-UCB score ensures that stale estimates (large $t_i$) contribute less to the score. After $\tau = \mathcal{O}(1/\lambda)$ steps without playing arm $i$, the weighted weakness decays to $\mathcal{O}(e^{-1})$, while the urgency term $\gamma(1 - e^{-\lambda \tau}) \approx \gamma(1 - e^{-1})$ provides a floor that prevents arm starvation, guaranteeing each arm is played $\Omega(T/K)$ times.

\textbf{Step 2: Per-regime UCB analysis.}
Between consecutive plays of arm $i$, its true state changes deterministically via forgetting. Each ``regime'' (interval between successive plays) lasts at most $\mathcal{O}(K)$ steps since urgency forces revisiting. Within each regime, the standard UCB bound of~\citet{auer2002finite} applies: the expected number of suboptimal plays is $\mathcal{O}(\log T / \Delta_i^2)$, contributing $\mathcal{O}(\Delta_i \cdot \log T / \Delta_i^2) = \mathcal{O}(\log T / \Delta_i)$ regret per arm.

\textbf{Step 3: Minimax rate.}
Summing the per-arm regret across $K$ arms:
\[
R_T \;\leq\; \sum_{i=1}^{K} \mathcal{O}\!\left(\frac{\log T}{\Delta_i}\right).
\]
The minimax regret is obtained by evaluating at the worst-case gap $\Delta_i = \Theta(\sqrt{K\log T / T})$ (all arms equally suboptimal at the hardest scale).
At this scale, $\log T / \Delta_i = \log T / \Theta(\sqrt{K\log T/T}) = \Theta(\sqrt{T\log T / K})$.
Summing over $K$ arms:
\[
R_T \;\leq\; K \cdot \mathcal{O}\!\left(\sqrt{\frac{T\log T}{K}}\right) = \mathcal{O}\!\left(\sqrt{KT\log T}\right).
\]

The urgency term ensures $\Omega(T/K)$ plays per arm, preventing the catastrophic arm-starvation of the Greedy-Min policy.
\end{proof}

\begin{proof}[Proof of Theorem~\ref{thm:lower_bound}]
We construct a hard family of forgetting-bandit instances that reduces to a standard $K$-armed bandit.

\textbf{Hard instance family.}
Consider forgetting-bandit instances with $\alpha = \beta = \lambda = 0$ (no learning, no forgetting) and $K$ arms with fixed initial knowledge levels $k_1, \ldots, k_K \in [0,1]$.
With all dynamics disabled, arm $i$ yields reward $k_i$ at every timestep regardless of whether it is played---identical to a standard $K$-armed stochastic Bernoulli bandit where arm $i$ has mean $k_i$.
This is a well-defined subfamily of the forgetting bandit model.

\textbf{Lower bound via reduction.}
By Theorem~15.2 of~\citet{lattimore2020bandit}, the minimax regret for the $K$-armed stochastic bandit satisfies:
\[
\inf_\pi \sup_\nu R_T(\pi, \nu) \;\geq\; \frac{1}{20}\sqrt{KT},
\]
where the supremum is over all bandit instances with arms in $[0,1]$.
Since the hard instance family above is a valid subset of forgetting-bandit instances, and every algorithm for the forgetting bandit is in particular an algorithm for this subfamily, the same minimax lower bound applies.

\textbf{Extension to all $\lambda \geq 0$.}
Any forgetting-bandit algorithm must perform well on the $\alpha{=}\beta{=}\lambda{=}0$ subfamily (since these are valid problem instances), so the $\Omega(\sqrt{KT})$ lower bound applies for all $\lambda \geq 0$.

Together with Theorem~\ref{thm:fucb}, F-UCB achieves $\mathcal{O}(\sqrt{KT\log T})$ regret, which is within a $\sqrt{\log T}$ factor of this lower bound. Since $\sqrt{\log T}$ grows sub-polynomially, F-UCB is near-minimax-optimal for the forgetting bandit.
\end{proof}
